\section{limitations \& future directions}
\label{sec:limitations}
We identify \sotopia as the first platform for a general and realistic evaluation of social intelligence in AI agents.
To better understand the social intelligence of AI agents, we discuss some future directions for \sotopia and the field of AI social intelligence.

\paragraph{Limitations of the simplified simulated ``world''} 
As every simulation is a simplification of the real world, \sotopia identifies several key components of realistic social interactions, while abstracting aspects of the real world.
First, we consider five types of social relationships in \sotopia. Future work could expand the type and granularity of social relationships (e.g., colleagues, classmates, etc.) in \sotopia. Different types of relationships would require agents to exhibit different social behaviors \citep{Jenkins2018-uw}, making the expansion of relationship types an important future research direction. 
Second, future work could expand the breadth of the character and social scenario pool in \sotopia to cover more social behaviors.
Third, \sotopia constrains the fixed turn-taking interaction to the \textit{dyadic context}, studying interactions between two agents. Future works could tackle more complex social interactions, such as multi-party interactions and those involving complex dynamics (e.g. asynchronous interactions, interruptions).



\paragraph{Social impact and ethical considerations}
Attributing human characteristics to AI systems risks anthropomorphizing them, which could lead to unrealistic expectations of AI systems, potential manipulation, and negative influence \citep{deshpande2023anthropomorphization}.
AI agents in \sotopia are not dedicated to a consistent human identity but rather role-play various characters across different scenarios.
This role-playing setting discourages AI systems with consistent human personalities, which could lead to anthropomorphism \citep{shanahan2023roleplay}.
The main goal of \sotopia is to evaluate the social intelligence of AI agents, and we do not intend to create AI agents that are indistinguishable from humans.
We consider the interactions that happened in \sotopia as simulacra of human interactions and such simulated interactions could help us better understand the social intelligence of AI agents, and explore various social phenomena \citep{Park2023GenerativeAI}.

Potential social stereotypes that are embedded in the automated evaluation system in \sotopia, as it is majorly supported by GPT-4 \citep{cheng-etal-2023-marked}.
Future work could investigate when such biases emerge, how they affect the evaluation, and how to mitigate them.
Identifying potential biases in \sotopia could also help scientists better understand social biases in the real world \citep{zhou-etal-2020-debiasing}.
Future work could also extend the evaluator with other systems, for example, Delphi \citep{jiang2022machines}.
Mitigating biases and stereotypes in interactive \sotopia-like systems could support the development of social AI agents that are more fair and inclusive.

Meanwhile, models learn to persuade or negotiate with humans, which may lead to social manipulation.
We do not endorse the use of \sotopia to create manipulative agents and will release \sotopia under the AI2 impact license\footnote{https://allenai.org/impact-license} to prevent misuse. 
Future work could further investigate the potential risks of AI anthropomorphism and manipulation and design more robust evaluation systems to mitigate these risks with \sotopia.




\paragraph{Improving LLM social intelligence}
Our \sotopia environment and  \sotopiaeval framework provide the opportunity for researchers to train more socially intelligent language agents.
As shown in section~\ref{sec:human_eval}, GPT-4 is able to provide reasonable evaluations for social interactions even for interactions involving humans. 
Future work could explore using the automated evaluation system to provide rewards to train LLMs with enhanced social intelligence.
